\begin{table}[h]
\centering
\small
\caption{Efeito de cada esquema de perturba\c{c}\~ao sobre o \emph{ground truth}: raz\~ao entre o GT recalculado a partir do invent\'ario perturbado e o GT original publicado, por sistema e escopo geogr\'afico.}
\label{tab:sm_stress_gt_ratio}
\begin{tabular}{clccc}
\toprule
& & \multicolumn{3}{c}{\textbf{GT novo / GT original}} \\
\cmidrule(lr){3-5}
\textbf{Sis.} & \textbf{Esquema} & \textbf{Global} & \textbf{Brasil} & \textbf{Estado} \\
\midrule
S1 & P1 & 2,000 & 2,000 & 2,000 \\
S1 & P2 & 0,373 & 0,296 & 0,288 \\
S1 & P3 & 1,282 & 1,328 & 1,332 \\
\midrule
S2 & P1 & 2,000 & 2,000 & 2,000 \\
S2 & P2 & 0,570 & 0,210 & 0,317 \\
S2 & P3 & 1,177 & 1,388 & 1,323 \\
\bottomrule
\end{tabular}
\caption*{\footnotesize Nota: raz\~ao $1{,}000$ indicaria aus\^encia de deslocamento. O P1 desloca proporcionalmente e de forma id\^entica nos tr\^es escopos; o P2 desloca de forma distinta por escopo, porque o fator de caracteriza\c{c}\~ao da am\^onia \'e regionalizado; o P3 desloca sem padr\~ao. O menor deslocamento observado entre as 18 c\'elulas \'e de 17,7\%, acima do limiar de $\pm 10\%$ usado na classifica\c{c}\~ao, de modo que nenhuma chamada pode satisfazer simultaneamente os crit\'erios de \emph{c\'alculo} e de \emph{memoriza\c{c}\~ao}. Fonte: autores.}
\end{table}